\chapter[1950 --- Kendall et al.]{1950: Edward Calvin Kendall, Tadeus Reichstein, and Philip Showalter Hench}
\label{ch:1950_Kendall}

\section*{Main Contribution}

\textit{Isolated and characterized the hormones of the adrenal cortex (cortisone and related steroids), revolutionizing the treatment of rheumatic and inflammatory diseases.}

\section{Official Citation}

for their discoveries relating to the hormones of the adrenal cortex, their structure and biological effects

\section{Background and Historical Context}

\subsection{Edward Calvin Kendall}

Edward Calvin Kendall was born on March 8, 1886, in South Norwalk, Connecticut, USA. He attended Columbia University, where he earned his Ph.D. in chemistry in 1912 under the direction of Charles Frederick Chandler. After a brief period of postdoctoral work, Kendall joined the staff of the Mayo Foundation for Medical Education and Research in Rochester, Minnesota, where he would spend the next four decades conducting research on the biochemistry of hormones. At the Mayo Clinic, Kendall had access to large quantities of adrenal gland tissue from cattle slaughtered for food, which provided the raw material for his isolation of adrenal cortex hormones.

\subsection{Tadeus Reichstein}

Tadeus Reichstein was born on July 20, 1897, in Włocławek, Poland (then part of the Russian Empire). His family emigrated to Switzerland when he was a child, and he was educated in Zurich, where he studied chemistry at the Swiss Federal Institute of Technology (ETH). He received his Ph.D. in 1921 and subsequently joined the Department of Pharmaceutical Chemistry at the University of Basel, where he developed an interest in the chemistry of natural products, particularly the hormones produced by the adrenal cortex.

Reichstein's contribution to the study of adrenal cortex hormones was primarily chemical in nature. He developed innovative techniques for the extraction, purification, and identification of steroid compounds, and he isolated several important hormones from adrenal cortex tissue, including cortisone (which he called ``substance E''). His work was characterized by remarkable precision and ingenuity, and his isolation of cortisone in pure crystalline form was a crucial step in the development of cortisone as a therapeutic agent.

\subsection{Philip Showalter Hench}

Philip Showalter Hench was born on February 28, 1896, in Pittsburgh, Pennsylvania, USA. He attended Lafayette College and the University of Pittsburgh, where he received his medical degree in 1920. After completing his training, he joined the staff of the Mayo Clinic as a rheumatologist---a specialist in diseases of the joints and connective tissues. It was at the Mayo Clinic that Hench would make his primary contribution: the demonstration that cortisone was a significant treatment for rheumatoid arthritis.

The scientific context of the 1930s and 1940s was one of intense interest in the adrenal cortex and its hormones. The adrenal glands---small endocrine glands located on top of the kidneys---produce a variety of hormones that regulate metabolism, immune function, and the body's response to stress. The hormones of the adrenal cortex, which include cortisol, aldosterone, and their precursors, are steroids---compounds derived from cholesterol. The isolation and characterization of these hormones was a major goal of biochemical research, and the clinical use of adrenal cortex hormones for the treatment of disease was an even more ambitious objective.

\section{The Discovery/Contribution}

The Nobel Prize was awarded to Kendall, Reichstein, and Hench jointly ``for their discoveries relating to the hormones of the adrenal cortex, their structure and biological effects.'' Their combined contributions---Kendall's isolation and characterization of cortisol, Reichstein's chemical analysis of adrenal cortex steroids, and Hench's demonstration that cortisone could cure rheumatoid arthritis---represent one of the great achievements in the history of endocrinology and pharmacology.

\subsection{Kendall's Isolation of Cortisol}

Edward Kendall's primary contribution was the isolation and characterization of cortisol (originally called ``compound F''), the principal glucocorticoid hormone produced by the adrenal cortex. Working at the Mayo Foundation, Kendall developed techniques for extracting and purifying adrenal cortex hormones from beef adrenal glands. In 1934, he isolated what he called ``compound E''---later identified as cortisone---and in 1937, he isolated ``compound F''---cortisol---and determined its chemical structure.

Kendall's isolation of cortisol was a landmark achievement in biochemistry. Cortisol is the body's primary stress hormone, and it plays essential roles in the regulation of metabolism, immune function, inflammation, and the body's response to physical and emotional stress. The ability to produce cortisol in pure form opened up entirely new possibilities for both research and clinical medicine.

\subsection{Reichstein's Chemical Analysis}

Tadeus Reichstein's contribution was the systematic chemical analysis of adrenal cortex steroids. Using innovative chromatographic techniques, Reichstein identified and characterized over 30 different steroid compounds produced by the adrenal cortex, including precursors and metabolites of cortisol and aldosterone. He also developed methods for the partial synthesis of cortisone from naturally occurring steroids, which would later be exploited for the large-scale production of the drug.

Reichstein's work was essential for understanding the biosynthetic pathways by which the adrenal cortex produces its hormones. He showed that the different adrenal cortex steroids are synthesized from cholesterol through a series of enzymatic steps, and that the relative amounts of each steroid are regulated by the pituitary gland through the hormone ACTH (adrenocorticotropic hormone). This understanding was crucial for the development of therapies based on adrenal cortex hormones.

\subsection{Hench and the Treatment of Rheumatoid Arthritis}

Philip Hench's contribution was the demonstration that cortisone was a significant treatment for rheumatoid arthritis and other inflammatory diseases. Hench had long been interested in the relationship between adrenal cortex function and arthritis, and he had observed that pregnant women and patients with jaundice often experienced temporary improvement in their arthritis symptoms. He hypothesized that these improvements might be related to changes in the body's production of adrenal cortex hormones.

In 1941, Kendall provided Hench with a small supply of compound E (cortisone), and Hench began testing it in patients with rheumatoid arthritis. The first patient, a 29-year-old woman named Anne, had been severely incapacitated by her disease and was unable to walk. After receiving cortisone, she showed dramatic improvement within days, regaining the ability to walk and perform daily activities. The results were so spectacular that Hench initially hesitated to publish them, fearing that they would not be believed.

Hench published his landmark paper in 1948, reporting the results of cortisone treatment in 14 patients with rheumatoid arthritis. The paper caused a sensation in the medical world, as it demonstrated for the first time that a chemical compound could produce dramatic improvement in a chronic inflammatory disease. Within a few years, cortisone was being used worldwide for the treatment of rheumatoid arthritis and other inflammatory conditions.

\section{Impact on Modern Medicine}

The discoveries of Kendall, Reichstein, and Hench have had a transformative impact on modern medicine, leading to the development of a vast array of therapies based on corticosteroid hormones.

\subsection{Corticosteroid Therapy}

The introduction of cortisone and its synthetic derivatives (prednisone, prednisolone, dexamethasone, and others) revolutionized the treatment of inflammatory and autoimmune diseases. Corticosteroids are now used to treat a wide range of conditions, including rheumatoid arthritis, lupus, inflammatory bowel disease, asthma, allergic reactions, and organ transplant rejection. They are among the most widely prescribed drugs in the world, and they have saved millions of lives.

The mechanism of action of corticosteroids was elucidated in the decades following their introduction. Corticosteroids work by binding to intracellular receptors that act as transcription factors, modifying the expression of genes involved in inflammation and immune function. By suppressing the production of inflammatory mediators such as prostaglandins and leukotrienes, and by inhibiting the activity of immune cells such as T lymphocytes and macrophages, corticosteroids reduce inflammation and suppress the immune response.

\subsection{Adrenal Insufficiency and Hormone Replacement}

The understanding of adrenal cortex function that emerged from the work of Kendall, Reichstein, and Hench has been essential for the diagnosis and treatment of adrenal insufficiency (Addison's disease), a condition in which the adrenal glands do not produce adequate amounts of cortisol and aldosterone. Patients with Addison's disease require lifelong replacement therapy with corticosteroid hormones, a treatment that is based directly on the discoveries of the 1950 Nobel laureates.

\subsection{Topical Corticosteroids}

The development of topical corticosteroid formulations---creams, ointments, and lotions that can be applied directly to the skin---has been one of the major advances in dermatology. Topical corticosteroids are used to treat a wide range of inflammatory skin conditions, including eczema, psoriasis, and allergic dermatitis. The development of these formulations, which deliver the anti-inflammatory effects of corticosteroids to the skin while minimizing systemic side effects, represents a direct application of the pharmacological principles established by Kendall, Reichstein, and Hench.

\subsection{Understanding and Managing Side Effects}

The widespread use of corticosteroids has also highlighted the importance of understanding and managing the side effects of these powerful drugs. Long-term use of corticosteroids can cause a range of adverse effects, including osteoporosis, diabetes, weight gain, cataracts, and immune suppression. The development of strategies to minimize these side effects---including the use of the lowest effective dose, the use of steroid-sparing agents, and the development of newer drugs with more selective mechanisms of action---has been an important area of pharmacological research that builds directly on the work of the 1950 Nobel laureates.

\subsection{Modern Glucocorticoid Research}

The understanding of glucocorticoid biology that emerged from the work of Kendall, Reichstein, and Hench has continued to evolve and expand. Research into the molecular mechanisms of glucocorticoid action has revealed that these hormones have complex and often tissue-specific effects, influencing the expression of thousands of genes in different cell types. This understanding has led to the development of more selective glucocorticoid receptor modulators (SEGRMs) that aim to retain the anti-inflammatory benefits of corticosteroids while reducing their metabolic and endocrine side effects.

\section{Legacy}

The legacies of Edward Kendall, Tadeus Reichstein, and Philip Hench are inextricably linked through their joint contribution to the discovery and clinical application of adrenal cortex hormones. Their work transformed the treatment of inflammatory and autoimmune diseases, established the field of corticosteroid pharmacology, and saved millions of lives around the world.

Kendall continued his research at the Mayo Clinic until his retirement, making additional contributions to the understanding of thyroid hormones and other endocrine systems. He received numerous honors, including the Presidential Medal of Science, and died on May 4, 1972, at the age of 86. Reichstein remained at the University of Basel for his entire career, continuing his research on steroid chemistry and natural products. He received numerous honorary degrees and awards and died on August 1, 1996, at the age of 99. Hench continued his research at the Mayo Clinic and became an outspoken advocate for the responsible use of corticosteroids. He died on March 30, 1965, at the age of 69.

The work of the 1950 Nobel laureates has had an enduring impact on medicine, and their legacy continues to shape the practice of rheumatology, endocrinology, and pharmacology. As new glucocorticoid-based therapies are developed and as our understanding of the molecular mechanisms of inflammation continues to deepen, the foundational work of Kendall, Reichstein, and Hench remains an essential reference point and an enduring source of inspiration.


\section{Modern Developments}

Kendall, Reichstein, and Hench's adrenal hormone work has evolved into modern endocrinology and immunology. Corticosteroids remain widely used for inflammatory and autoimmune diseases, though with awareness of long-term side effects. Selective glucocorticoid receptor modulators (SEGRMs) aim to retain anti-inflammatory effects while minimizing metabolic side effects. Understanding of adrenal insufficiency has led to better hormone replacement therapies. Cushing's syndrome treatment now includes targeted therapies (osilodrostat, pasireotide).
